\documentclass[sigconf,anonymous]{aamas}
\usepackage{amsmath}
\usepackage{booktabs}
\usepackage{array}
\usepackage{url}
\input{aamas_metadata}
\vfuzz=1.5pt
\acmSubmissionID{Anonymous submission}
\submissionType{Research Paper Track}

\title[Earning Roles]{Earning Roles: Situated Peer Evaluation and Dynamic Responsibility Assignment Inspired by Human Social Collaboration}
\author{Anonymous Authors}
\affiliation{\institution{Anonymous Institution}}
\email{omitted}
\ccsdesc[500]{Computing methodologies~Multi-agent systems}
\ccsdesc[300]{Computing methodologies~Cooperation and coordination}
\keywords{multi-agent systems, human social collaboration, dynamic roles, peer evaluation, online adaptation, responsibility attribution}

\newcommand{\rare}{\textsc{rare}}
\newcommand{\unknown}{\textsc{unknown}}

\begin{abstract}
In human collaboration, expertise is rarely established by self-declaration alone. It emerges as partners use one another's work, evaluate concrete deliveries, and decide whom to entrust with similar responsibilities next. We translate this principle into \emph{earning roles}: an agent's role is its task-situated eligibility for future responsibility, not a permanent title or a context-free reputation score. The system records a recipient's acceptance, revision, or rejection of a versioned artifact together with its use context as typed, noisy, auditable experience. Peer evaluation is exposed to subsequent selection as an observation, while producer responsibility is decided by an independent attribution check. Selector parameters change only after a later assignment has been executed and its independent quality, responsibility, and complete cost have passed validation; valid feedback may record either success or failure, whereas missing or unattributable outcomes do not update the policy. We evaluate this complete chain in an ArtifactRole track and isolate local candidate sets, delayed feedback, and online service behavior in a complementary PeerSelect track. All comparisons use matched candidates, public information, budgets, and feedback timing. \rare{} attains a future quality--cost utility of xx, differing from the same-information primary baseline by xx (stream-level 95\% CI: xx); judgment Brier score, false producer-attribution rate, and update p95 are xx, xx, and xx. Ablations separate the effects of peer information, responsibility validation, and delayed updating.
\end{abstract}

\begin{document}
\maketitle

\section{Introduction}
In human societies, professional standing is seldom determined by self-description. It develops when collaborators use a person's work, evaluate whether the delivery was useful or required repair, and decide whether to entrust related responsibilities to that person again. We adopt this rule without treating social reputation as ground-truth ability: value is inferred from concrete work and tested by subsequent delegation.

Role prompts and fixed handoff orders in multi-agent workflows usually lack this feedback loop. A role name states what an agent is expected to do, not what it has demonstrated in collaboration. When tasks or member performance change, a workflow that ignores collaboration experience continues to reuse its original allocation. We therefore define \emph{earning roles} as acquiring eligibility for future responsibility in a task context. A role is not a permanent position; expertise is not self-declared; and earning does not confer ownership forever.

Our evaluator is the agent that actually receives and uses a versioned delivery. Its acceptance, revision, or rejection is stored with the task, artifact, and use trace as a \emph{collaboration experience}. A selector reads these experiences on a later task. An independent responsibility check determines which outcomes are attributable to the producer, and the selected agent's subsequent quality and complete cost determine whether that allocation should be repeated. Peer evaluation, responsibility attribution, and future assignment thus retain distinct meanings.

Consider a queue-processing workflow. Agent $B$ supplies a priority module and agent $A$ integrates it into a consumer. When another task requires queue processing, the selector may use $A$'s account of whether and how the module worked, rather than rely only on $B$'s self-description. Yet a rejection is not automatically evidence of a producer defect. The same failed team outcome can arise because $B$ violated the module contract or because $B$'s module was correct while $A$'s consumer was faulty. The system must retain the recipient judgment, actual use, and independent quality check to distinguish these cases.

We instantiate this principle as the earning-roles protocol. The selector combines current requirements, candidate information, and prior experiences before execution. Positive and negative recipient judgments remain observations; producer responsibility is validated separately. Publishing an observation can change a later selector input but cannot directly update its parameters. Only a later, executed assignment with independently validated quality and complete cost can update the selector. A judgment is therefore not converted into immediate reward or punishment.

We call the responsibility-qualified component \emph{responsibility-aware role evidence} (\rare{}) and store it separately from raw evaluation. Figure~\ref{fig:overview} shows how this distinction supports subsequent delegation. The central test compares \rare{} with acceptance-only, terminal-only, and same-information contextual selectors to measure its effect on unseen-task collaboration quality and its evaluation, validation, and update costs.

The paper makes three contributions.
\begin{itemize}
  \item We define task-situated collaboration experience that records who evaluated whose versioned delivery, how it was used, and which outcomes are attributable to the producer. Both positive and negative evaluations remain auditable in context.
  \item We separate observation publication from delayed policy updating. Published evaluations inform later choices, whereas parameters update only from independent outcomes of executed assignments; factorial ablations isolate the two paths.
  \item We introduce an ArtifactRole track for the complete delivery--evaluation--delegation--utility chain and a PeerSelect track for local selection, delay, drift, and service behavior. Same-information baselines and responsibility interventions separate information value, attribution error, and gains from additional information.
\end{itemize}

\begin{figure*}[t]
\centering
\includegraphics[width=\textwidth]{figures/overview.pdf}
\caption{From peer evaluation to future responsibility. (a) Both cases receive a revision judgment, but responsibility differs. An independent check $Q$ of producer $B$'s original version $v0$, together with the recipient's use trace, determines whether the observation forms evidence $e$ attributable to $B$. (b) Collaboration experiences are read before a later assignment; the executed assignment changes subsequent selection only after outcome validation. File changes, candidate state, and selections are schematic.}
\Description{Two queue-processing cases contrast a producer defect with recipient-owned integration work. Versioned evidence enters a later selector before assignment; only a validated later outcome updates future state.}
\label{fig:overview}
\end{figure*}

\section{Related Work and Positioning}
Peer feedback has recently been used to adapt multi-agent collaboration. Meta-Team collects post-task accounts of how teammates affect downstream execution and updates teammate profiles and collaboration arrangements \cite{hao2026metateam}. RAPS maintains peer reputation from first-hand content evaluations and witness reports \cite{li2026raps}. Sero studies contract-preserving role evolution \cite{ge2026sero}, C3 studies exact credit assignment among cooperative agents \cite{chen2026c3}, and DecisionBench delegates from externally judged capability cards built from prior trajectories \cite{gao2026decisionbench}. These works cover collaborative reflection, witness propagation, contract-bound roles, credit assignment, and historical capability summaries, respectively. They do not jointly bind a recipient's use of one versioned delivery, a producer-responsibility gate, and a later executed responsibility opportunity into one auditable chain. Our analysis separates the evaluation of a concrete delivery, producer attribution, and the utility of subsequent delegation.

Computational trust and reputation provide the broader foundation for using interaction history in open multi-agent systems \cite{sabater2002reputation,pinyol2013trust}. Partner-selection methods adapt from social-game payoffs \cite{leung2024partner,leung2025partner,russell2026partner}, while language-agent societies use shared assessments or repeated interaction to sustain cooperation \cite{ren2026reputation,piatti2024cooperate,akata2025repeated}. Their interaction-level signals do not expose ownership of a producer artifact inside a recipient's integration path.

Graph and role systems reorganize agents from traces, routing outcomes, contexts, or subtask performance \cite{yang2025agentnet,zhuge2024gptswarm,zhang2026dynamic,li2026morse}. General multi-agent frameworks likewise make role-based workflows practical \cite{wu2023autogen,metagpt2024,agentverse2024}. Our object is narrower: an action on a specific delivery is separated from producer responsibility and from the payoff of a later assignment. Table~\ref{tab:positioning} summarizes this boundary.

\begin{table*}[t]
\centering\scriptsize
\setlength{\tabcolsep}{2.7pt}
\begin{tabular}{p{.17\textwidth}p{.22\textwidth}p{.19\textwidth}p{.17\textwidth}p{.16\textwidth}}
\toprule
Method family & Observation and update & Responsibility / credit & Experimental object & Difference tested here \\
\midrule
Meta-Team; RAPS \cite{hao2026metateam,li2026raps} & Post-task accounts or first-/second-hand evaluation; profiles or reputation update & Teammate or content usefulness & Adaptive collaboration / publishing & Versioned use, independent producer gate, later assignment utility \\
Sero; C3; DecisionBench \cite{ge2026sero,chen2026c3,gao2026decisionbench} & Contract evolution, precise team credit, or judged trajectory cards & Role constraints, contribution credit, or capability summaries & Reasoning, cooperation, or delegation & Recipient action and producer attribution remain separate \\
Partner selection / reputation \cite{leung2025partner,russell2026partner,ren2026reputation} & Interaction outcomes and shared assessments & Peer payoff or pooled reputation & Social games / language-agent society & Producer--recipient artifact ownership \\
AgentNet; GPTSwarm \cite{yang2025agentnet,zhuge2024gptswarm} & Traces, graph structure, and routing outcomes & Route or role utility & Graph / workflow coordination & Action bound to one delivery and future owner \\
Dynamic roles / experts \cite{zhang2026dynamic,li2026morse} & Context and role/subtask performance & Role or expert utility & Debate / mixture-of-experts workflow & Responsibility gate and delayed correction \\
Finite-time bandit \cite{auer2002finite} & Selected-arm reward after execution & Arm-level reward & Pre-execution choice & Context, ownership, and delay require qualification \\
\bottomrule
\end{tabular}
\caption{Comparison dimensions. Adapted experiments preserve matched public information, feedback timing, and complete cost; rows denote different research objects rather than interchangeable implementations.}
\label{tab:positioning}
\end{table*}

Cooperative multi-agent learning separates shared return from individual contribution through counterfactual credit or value factorization \cite{foerster2018coma,rashid2020qmix}. These methods explain why team success is insufficient for individual credit, but their reward contracts do not identify ownership of a concrete artifact within recipient integration. Contextual bandits and delayed-feedback learning provide strong update controls \cite{li2010contextual,joulani2013delayed}; they do not supply the proposed event semantics. Thus, the relevant distinction is both temporal and semantic: source evaluation is visible before a later choice, whereas target quality exists only after execution and cannot be backfilled into the earlier decision.

\section{Problem Formulation}
\label{sec:problem}
\subsection{Tasks, assignments, and collaboration experience}
At task $t$, context $x_t$ and legal candidate set $C_t$ specify requirements and delivery responsibility. Before execution, the selector chooses
\begin{equation}
 a_t\sim\pi_{\psi,\Theta_t}(\cdot\mid x_t,C_t,R_t),\qquad
 p_t=\pi_{\psi,\Theta_t}(a_t\mid x_t,C_t,R_t),
\label{eq:choice}
\end{equation}
where $a_t\in C_t$, $p_t\in(0,1]$, $R_t$ is the public experience available at that time, $\psi$ denotes fixed scoring parameters, and $\Theta_t$ is online state. Public records and parameters are maintained separately, so a new evaluation can change the choice input without immediately changing parameters.

A source delivery forms the record
\begin{equation}
 e_i^{\mathrm{src}}=(x_i,C_i,a_i,o_i,j_i,J_i,m_i,Q_i,Y_i^{\mathrm{src}},\tau_i,c_i^{\mathrm{src}}).
\label{eq:event}
\end{equation}
Here $o_i$ is producer $a_i$'s versioned artifact; $j_i$ is the recipient that uses it; $J_i$ is sealed before the recipient acts; $m_i$ is the subsequent use, revision, or rejection; $Q_i$ checks the original delivery against the producer's responsibility; $Y_i^{\mathrm{src}}$ is the source episode's independently defined terminal or later-use outcome, or $\bot$ if the contract defines none; $\tau_i$ is publication order; and $c_i^{\mathrm{src}}$ records model, tool, token, and wall-clock cost. Source outcomes support observation and audit, not future-assignment credit.

For a target assignment $j$, the independent outcome is $Y_j^{\mathrm{tar}}$. After the assignment is sealed, its candidates and probability are consistent, and execution completes, $Y_j^{\mathrm{tar}}$ enters $\operatorname{validate\_later}$ with that assignment to produce feedback $L_j$. Only an $L_j$ that passes responsibility, quality, complete-cost, and idempotency checks updates $\Theta$. Hence $Y_i^{\mathrm{src}}$ and $Y_j^{\mathrm{tar}}$ are distinct, and neither is retroactively available to an earlier choice.

\subsection{Visible information and objective}
The selector reads a public projection $\tilde e_i=\Pi_{\mathrm{pub}}(e_i)$ rather than the audit record. The projection preserves agreed task fields, candidate identity, recipient judgment, actual action, and responsibility type; hidden answers, private tests, and future outcomes remain unavailable. Given context $x$, candidate $a$, and public history $R_t$, situated assignment value is
\begin{equation}
 V_t(a\mid x,R_t)=\mathbb{E}[U(Y_j^{\mathrm{tar}},c_j^{\mathrm{tar}})\mid x,R_t,a],
\label{eq:value}
\end{equation}
where $U$ is the frozen quality--complete-cost utility. This quantity asks where a candidate should again receive responsibility, not whether it has context-free global ability. Producer quality, recipient evaluation, and team outcome are recorded separately. Team success does not prove that every member performed well; a negative judgment does not automatically label the producer. Self-description supplies at most an initial prior and never verified evidence.

\section{Earning Roles}
\label{sec:method}
The protocol preserves evaluations from actual collaborators, organizes task-relevant history, assigns the next responsibility, and calibrates the selector from the executed result. Responsibility validation distinguishes what a peer observed from what can be attributed to the producer. The auditable order is
\begin{equation}
\text{delivery}\rightarrow\text{sealed judgment}\rightarrow\text{use}\rightarrow\text{publication}
\rightarrow\text{assignment}\rightarrow\text{execution}\rightarrow\text{validated update}.
\label{eq:order}
\end{equation}
A later result cannot change an earlier assignment.

\subsection{Observation and responsibility}
We define observation completeness $O_i\in\{0,1\}$ and producer-attribution eligibility $G_i\in\{0,1\}$. $O_i=1$ when the evaluation, versioned delivery, recipient action, and independent check are completely bound under the prescribed time order. It does not depend on whether the recipient accepts, revises, or rejects; all three remain collaboration experience.

$G_i=1$ adds responsibility requirements: the original delivery maps uniquely to a predeclared producer duty, the recipient actually used it, the independent check supports attribution, and the evidence does not mix recipient-owned integration changes. Responsibility follows the task contract, artifact lineage, and actual modifications rather than the evaluator's assertion. Recipient-only work, mixed ownership, or unresolved responsibility does not generate producer reward or penalty. Records with $O_i=1,G_i=0$ remain typed contextual observations but cannot be reinterpreted as producer-quality labels. In the queue example, $A$'s integration difficulty may affect future collaboration choice, but it cannot become evidence that $B$'s module is defective.

\subsection{History and assignment}
At read cut $w_t$, candidate version $v$ exposes
\begin{align}
\mathcal E_t^{\mathrm{obs}}(v)&=\operatorname{Resolve}_{w_t}
 \{\tilde e_i:O_i=1,a_i=v,\tau_i\leq w_t\},\\
\mathcal E_t^{\mathrm{cred}}(v)&=\operatorname{Resolve}_{w_t}
 \{\tilde e_i:O_i=1,G_i=1,a_i=v,\tau_i\leq w_t\}.
\label{eq:views}
\end{align}
$\operatorname{Resolve}_{w_t}$ deduplicates by source and retains only the valid version known at $w_t$. Late evaluations affect only later choices. Globally, $\mathcal R_t^{\mathrm{obs}}$ retains complete recipient observations, including responsibility-unknown records, whereas $\mathcal R_t^{\mathrm{cred}}$ contains attributable evidence. Publishing either view does not change $\Theta$, the update count, or the policy digest. Only a validated later $L_j$ can trigger parameter updating.

History is divided into an archive and a new delta:
\begin{align}
 P_{v,k}&=\operatorname{summarize}\{\tilde e_i\in\mathcal E_t^{\mathrm{obs}}(v):\tau_i<\kappa_k\},\\
 D_{v,k,t}&=\operatorname{summarize}\{\tilde e_i\in\mathcal E_t^{\mathrm{obs}}(v):\kappa_k\leq\tau_i\leq w_t\},\\
 z_t(v)&=[\phi(x_t)\,\|\,\phi(P_{v,k})\,\|\,\phi(D_{v,k,t})\,\|\,b_t(v)].
\label{eq:representation}
\end{align}
The structured summaries retain source identifiers; empty history is explicit. Corrections rebuild an affected archive before its next read, advance the archive version, and repartition the delta. The frozen encoder $\phi$ represents the task, archive, and delta; $b_t(v)$ contains public version and coverage fields.

A fixed head and candidate-specific residual define
\begin{equation}
 s_t(v)=g_\psi(z_t(v))+\theta_{v,t}^{\top}\phi(x_t).
\end{equation}
The first term reads current collaboration experience; the second calibrates scores from previously validated execution outcomes. With temperature $T>0$ and exploration $0<\varepsilon<1$,
\begin{align}
 q_t(v)&=\frac{\exp(s_t(v)/T)}{\sum_{u\in C_t}\exp(s_t(u)/T)},\\
 \pi_t(v)&=(1-\varepsilon)q_t(v)+\varepsilon/|C_t|,\qquad a_t\sim\pi_t.
\label{eq:policy}
\end{align}
Every candidate therefore retains nonzero opportunity. The assignment record seals the candidate set, sample, probability, and history version before target execution. Unselected candidates receive no outcome labels. Complete cost includes selected-agent work, recipient integration, evaluation, and validation.

\begin{figure*}[t]
\centering
\includegraphics[width=\textwidth]{figures/method_state.pdf}
\caption{Experience reading and delayed updating. (a) Effective history is split into archive $P$ and delta $D$, then combined with task representation and public metadata. The candidate and propensity are sealed before execution. (b) After execution, the matching independent quality and cost outcome must pass identity, responsibility, and once-only checks before updating the residual used by later decisions.}
\Description{The upper path reads a versioned archive, delta, task, and public metadata to sample a candidate before execution. The lower path validates the executed assignment and updates only future residual state.}
\label{fig:method-state}
\end{figure*}

\subsection{Delayed updating}
Publication and parameter updating are independent:
\begin{align}
 \operatorname{publish}(\tilde e_i)&:R\leftarrow R\cup\{\tilde e_i\},\quad
 \Theta,\ \text{update count, policy digest unchanged},\\
 \operatorname{update}(L_j)&:\Theta^{+}=U_{\mathrm{delay}}(\Theta;L_j).
\label{eq:update}
\end{align}
$L_j$ binds the pre-execution assignment, propensity, and input snapshot to the independent target outcome, responsibility check, quality, and complete cost. The frozen activation manifest specifies representation, updater, initialization, clipping, batching, and stopping constants. Validation checks assignment order, candidate identity, outcome completeness, reward semantics, and once-only use. Valid successes and failures both form feedback; missing or unattributable labels do not update. Corrections and delayed records affect only choices after their arrival.

\section{Benchmarks and Baselines}
\paragraph{ArtifactRole.} The primary track uses TeamBench's dependency structure \cite{kim2026teambench} to expose producer delivery, recipient use and evaluation, responsibility checks, later assignment, and complete cost. DIST1 queue-race and PIPE3 stream-processing separate development and confirmation roots; changing a random seed does not create a new structural root. Producer contract conformance and terminal usability are scored separately because a recipient may normalize a nonconforming delivery into a usable result. Every executed arm, including no-update, receives terminal evaluation. Valid assignment credit is required for updating, not for entry into the intention-to-treat outcome denominator.

\paragraph{PeerSelect.} The secondary track uses a pinned local graph-interaction benchmark to isolate selected-only payoff, local candidate sets, online updating, drift, and service cost. It cannot establish artifact-specific judgment or producer attribution. The two tracks share event and selector interfaces while retaining separate tasks, labels, scorers, denominators, and claims.

Agent benchmarks supply complementary coordination and interaction coverage \cite{zhu2025multiagentbench,agashe2025llmcoordination,liu2024agentbench,appworld2024,yao2024taubench}. They do not, by default, expose the proposed producer--recipient lineage. All primary arms receive matched candidate menus, initial opportunities, model/API/tool budgets, exploration, feedback arrival schedules, failure handling, and complete-cost accounting.

\begin{table*}[t]
\centering\scriptsize
\setlength{\tabcolsep}{3pt}
\begin{tabular}{p{.16\textwidth}p{.42\textwidth}p{.33\textwidth}}
\toprule
Arm & Read / update contract & Alternative explanation tested \\
\midrule
Uniform & Candidate menu; fixed local-uniform sampling & Random allocation and exploration floor \\
No update & Initial prior and current context; feedback ignored & Static execution is sufficient \\
Raw acceptance & Recipient acceptance / rejection only & Responsibility filtering adds no value \\
Terminal only & Independent terminal outcome only & Situated evaluation adds no information \\
Contextual trust / bandit & Same context, evaluation, propensity, delay, and budget as \rare{} & Ordinary contextual learning explains the effect \\
Pooled controller & Permitted public history at matched total capacity & Shared control explains the gain \\
Closest published adapter & Public-information adapter preserving source semantics & Existing role / organization method is sufficient \\
\rare{} & Observation and responsibility separated; validated selected assignment updates & Complete evaluation--delegation--calibration mechanism \\
\bottomrule
\end{tabular}
\caption{Predeclared baselines. Same-information comparisons preserve context, feedback timing, and complete budget.}
\label{tab:baselines}
\end{table*}

RLS, online logistic regression, SGD, and periodic refitting are compared under the same selection interface. This separates updater computation from responsibility-aware information use. Continual-learning and multi-agent non-stationarity motivate the drift and forgetting analyses \cite{khetarpal2022continual,hernandezleal2017nonstationarity}; agent ranking and human--LLM teaming motivate generality and deployment boundaries \cite{lanctot2025soft,yurrita2026guidelines}.

\section{Experimental Questions and Matrix}
\textbf{RQ1 (information value).} Does eligible situated evaluation predict an independent producer-contract or later-use result beyond raw acceptance and terminal-only feedback? Measures are held-out log loss, Brier score, AUC when defined, calibration, and incremental value on an independent root.

\textbf{RQ2 (assignment causality).} Does public evidence change an assignment sealed before target execution, including transfer to a future owner different from the source producer? Measures are assignment-change probability, propensity-weighted future quality, adoption, rework, and complete cost. Evidence about $B$ cannot be consumed as evidence about $C$.

\textbf{RQ3 (online behavior).} Does delayed selected-only updating satisfy the service contract? Measures are publish, read/assignment, and update p50/p95 latency; service lag; backlog; state bytes; CPU/GPU/RAM; token/tool/API use; and replay time.

\textbf{RQ4 (stability and boundaries).} How does the method behave under task or peer drift, delayed or out-of-order feedback, version replacement, recipient-only work, mixed ownership, scorer failure, and resource failure? Measures are peak and mean old-root forgetting, recovery window, \unknown{} rate, correction consistency, and false producer attribution.

The uncertainty unit is an independently reset stream; events within a stream are dependent. We use paired stream-level estimates, bootstrap or randomization intervals over streams, and preregistered multiplicity control. Missing, invalid, and unknown events remain in their declared denominators. H1 is judgment Brier score, H2 future quality--complete-cost utility, H3 update p95 and backlog, and H4 false producer-attribution rate. H2 is higher-is-better; the others are lower-is-better. Comparison direction, precision target, 95\% interval, and stopping rule are sealed before confirmation outcomes are read.

\begin{table*}[t]
\centering\scriptsize
\setlength{\tabcolsep}{2.6pt}
\begin{tabular}{p{.07\textwidth}p{.16\textwidth}p{.25\textwidth}p{.27\textwidth}p{.15\textwidth}}
\toprule
Cell & Track / root & Strategy / intervention & Matched comparison & Measurement / result \\
\midrule
RQ1-I & ArtifactRole / dev & Raw vs. terminal vs. \rare{} & Same delivery, recipient observation, legal feedback & Brier/log loss, calibration: xx \\
RQ1-II & ArtifactRole / confirm & Contextual trust vs. \rare{} & Same read cut, delay, propensity, model, budget & Incremental prediction, transfer: xx \\
RQ2-I & ArtifactRole / confirm & No evidence vs. public evidence & Preview and assignment before target execution & Change, quality, adoption, rework: xx \\
RQ2-II & ArtifactRole / confirm & Delayed-only vs. evidence+delay & Same target and selected-only outcome & Future utility, complete cost: xx \\
RQ2-III & ArtifactRole / confirm & \rare{} vs. pooled controller & Same capacity and candidate menu & Utility, concentration, propensity: xx \\
RQ3-I & PeerSelect / dev stream & Uniform, no/history, updater & Pinned graph and call/token budget & Payoff, p50/p95, state bytes: xx \\
RQ3-II & PeerSelect / drift & Update with/without replay & Fixed peer/task drift and delay & Recovery, forgetting, backlog: xx \\
RQ4-I & ArtifactRole / eligible roots & Responsibility mutations & Producer, recipient, mixed, unknown, failure & Attribution, UNKNOWN, no-op: xx \\
RQ4-II & ArtifactRole / eligible roots & Remove one mechanism & Gate, evidence, delay, locality, updater & Utility, safety, cost, interaction: xx \\
\bottomrule
\end{tabular}
\caption{Preregistered matrix. Each cell is bound to a structural root, split, policy, seed, budget, scorer, cost record, and raw-event manifest.}
\label{tab:matrix}
\end{table*}

We distinguish $N_{\mathrm{scheduled}}$ eligible targets, $N_{\mathrm{started}}$ targets actually started after selection, $N_{\mathrm{outcome}}$ all started targets including failures and \unknown{}/invalid events, $N_{\mathrm{credit}}$ the selected-only and once-only validated subset, $N_{\mathrm{obs}}$ complete source observations, and $N_{\mathrm{attr}}$ attributable source evidence. Always $N_{\mathrm{credit}}\leq N_{\mathrm{outcome}}$; credit eligibility never filters the H2 intention-to-treat denominator. Predictive metrics use only their predefined scorable-label subset while reporting the full denominator and reasons for missingness.

Each request reserves attempts, prompt and completion tokens, tool time, and storage before dispatch. Observed costs from errors are retained. Complete cost covers producer execution, recipient integration, judgment service, communication, retries, replay, validation, and returned artifacts. We record configuration, model route/version, code/data hashes, seed, hardware, timestamps, raw responses, token use, actions, failures, unknown reasons, and per-sample results. Reports include stream-level means, medians, 95\% intervals, effect sizes, root heterogeneity, ITT and per-protocol denominators, and quality--cost Pareto points.

\paragraph{Mechanism ablations.} Count-only and judgment-masked controls preserve handoff count, arrival time, and cost while removing judgment category. The $2\times2$ design crosses public evaluation visibility $e$ with validated delayed updating $u$. With expected utility $V_{eu}$,
\begin{equation}
 \Delta_{\mathrm{obs}|\mathrm{upd}}=V_{11}-V_{01},\quad
 \Delta_{\mathrm{upd}|\mathrm{obs}}=V_{11}-V_{10},\quad
 I=V_{11}-V_{10}-V_{01}+V_{00}.
\label{eq:factorial}
\end{equation}
These contrasts distinguish information, updating, and interaction on the same utility scale. Responsibility interventions fix the producer delivery and vary recipient-owned integration; version tests cover duplicates, late records, and corrections across archive cuts.

\section{Results}
\begin{table}[t]
\centering\scriptsize
\setlength{\tabcolsep}{2.2pt}
\begin{tabular}{p{.39\columnwidth}p{.16\columnwidth}p{.18\columnwidth}p{.18\columnwidth}}
\toprule
Endpoint & \rare{} & Same-info. baseline & Difference / CI \\
\midrule
Future quality & xx & xx & xx \\
Complete cost & xx & xx & xx \\
Quality--cost utility & xx & xx & xx \\
Assignment change & xx & xx & xx \\
Judgment Brier & xx & xx & xx \\
Update p95 / state bytes & xx & xx & xx \\
Old-root forgetting & xx & xx & xx \\
UNKNOWN / false attribution & xx & xx & xx \\
\bottomrule
\end{tabular}
\caption{Primary results; intervals are computed over independently reset streams.}
\label{tab:results}
\end{table}

\paragraph{RQ1: information value.} In ArtifactRole, \rare{} differed from raw acceptance, terminal-only feedback, and same-information contextual trust in Brier score by xx, xx, and xx (stream-level 95\% CIs: xx, xx, and xx). Log loss, calibration error, and cross-root transfer were xx, xx, and xx. The evaluation included xx valid observations, an \unknown{} rate of xx, and judgment-service cost of xx. Predictive metrics remain distinct from team utility.

\paragraph{RQ2: assignment change.} Candidate sets, public history, propensities, and samples were sealed before target execution. Under \rare{}, assignment-change rate, future quality, rework, and quality--complete-cost utility were xx, xx, xx, and xx (95\% CI: xx). All scheduled and started targets entered the H2 ITT denominator, including quality-check, \unknown{}, invalid, resource, and scorer failures. Only the subset passing \texttt{policy\_update\_allowed} produced $L_j$ for delayed update and assignment credit.

\paragraph{RQ3: service behavior.} In PeerSelect, publish, read/assignment, and update p50/p95 latency were xx, xx, and xx. Backlog, state bytes, model calls, tokens, replay time, and complete service cost were xx, xx, xx, xx, xx, and xx; the difference from matched-budget updaters was xx. This track establishes the online service boundary but does not substitute for ArtifactRole attribution evidence.

\paragraph{RQ4: responsibility and drift.} Under recipient-only repair, mixed ownership, unknown responsibility, version replacement, out-of-order feedback, and resource failure, false producer attribution, \unknown{} precision, observation publication, valid update, drift recovery, old-root forgetting, and no-op rate were xx, xx, xx, xx, xx, xx, and xx (95\% CI: xx). Records whose responsibility could not be independently determined remained contextual observations and did not trigger producer reward or an ineligible parameter update.

\begin{table}[t]
\centering\scriptsize
\setlength{\tabcolsep}{3pt}
\begin{tabular}{llcc}
\toprule
Evaluation & Update & Future utility (95\% CI) & Cost \\
\midrule
Masked & Off & xx & xx \\
Visible & Off & xx & xx \\
Masked & On & xx & xx \\
Visible & On & xx & xx \\
\bottomrule
\end{tabular}
\caption{Evaluation--update factorial ablation.}
\label{tab:factorial-results}
\end{table}

All interpretations passed the corresponding structural-root, candidate-menu, public-information, timing, budget, and stopping-rule checks. ITT and per-protocol denominators, costs of failed runs, unknown reasons, and version corrections are reported together. Responsibility safety, prediction accuracy, and team utility are not treated as substitutes.

\section{Discussion and Limitations}
The method applies when responsibility boundaries, artifact versions, and later outcomes can be recorded explicitly. It does not discover open-ended workflows or infer general expertise from aggregate success. PeerSelect describes responsibility opportunities within local candidate sets; broader organizational restructuring requires candidate discovery and a new cost contract.

Human-inspired collaboration does not make peer evaluation ground truth. A recipient may misjudge, report selectively, or attribute its own integration difficulty to the producer. Complete but unattributable records remain contextual observations rather than producer rewards, and missing or invalid records do not update. This protects attribution at the cost of fewer learning examples. Cross-recipient calibration and adversarial-feedback modeling remain outside the present responsibility boundary.

Waiting for an executed future assignment grounds updating in actual allocation outcomes, but delays adaptation to change. Waiting alone does not establish attribution; independent responsibility checks remain necessary. The bounded selected-only update does not solve multi-step credit assignment over long chains. Frozen representations and a small update head reduce service cost but may under-express rapid semantic or relational drift; updating a larger encoder instead raises latency and forgetting risk. Real-time adaptation must therefore be evaluated jointly through new-task adaptation, old-task retention, and complete service cost.

Finally, complete lineage and audit records impose storage, privacy, and communication costs. Compressed delivery summaries reduce cost but may discard details needed for judgment; cross-tool or cross-organization deployments may not share all responsibility fields. Practical systems must trade replayability against data minimization and preserve a safe unknown branch for absent information.

\section{Conclusion}
Earning roles transfers a principle of human collaboration to multi-agent systems: future responsibility should depend on how partners evaluate concrete work and on whether subsequent delegation succeeds. Recipient evaluation forms task-situated role evidence; independent validation separates observation from producer attribution; and later quality and complete cost calibrate the selector. The resulting object is directly testable: whether collaborator evaluation changes future assignments and improves team performance on new tasks without sacrificing responsibility safety, service behavior, stability, or cost.

\balance
\clearpage
\bibliographystyle{ACM-Reference-Format}
\bibliography{references}
\end{document}
